\scanpage{110}

pour le changement de base)

\unsure{ens|en|ens.} des couples $(u,v)$ de \unsure{sections|section} \inserted{(mod scalaires inv.)} \unsure{abs$^{\mathrm{mt}}$} nilpotentes $\mathcal M$ \deleted{\unsure{tels q}} (i.e. de tr et dét nuls), \deleted{\unsure{distincts}}\inserted{non prop.} fibre par fibre $\simeq$ ens des réflexions de $\mathcal M$ $(\operatorname{tr}=0,\det=-1)$

\deleted{Mais les réflexions correspondent}

Si on veut passer aux clôtures projectives, on note que (comme car. rés $\ne2$) on a

\[
\mathcal M\simeq\underset{\substack{\Vert\\\tilde{\mathfrak g}}}{\unsure{\gamma|\mathfrak g}\mathcal M}+\unsure{\underline O_S|\mathcal O_S}
\]

$\hat\Sigma_{\mathcal M}\subset\deleted{\unsure{\text{variété}}}\check{\mathbb P}(\mathcal M)$ défini par l'équation \inserted{(dans $\check{\mathbb V}(\gamma\mathcal M\oplus\mathcal O_S)$)}

\[
\det(u)+\deleted{\unsure{\lambda}}t^2=0\qquad\text{i.e.}\quad\det(u+t\operatorname{id})=0
\]

Ou encore, en posant $v=u+t\operatorname{id}$

\[
\det v=0
\]

i.e. \deleted{le cône}\inserted{c'est la} quadrique projective définie par $\det$ de $\mathcal M$. Or on sait qu'on a un iso. canonique \deleted{avec}\inserted{entre celle-ci et} $X_{\mathcal M}\times X_{\mathcal M}$, \inserted{$F_{\mathcal M}:X_{\mathcal M}\times X_{\mathcal M}\xrightarrow{\sim}\hat\Sigma_{\mathcal M}$,} qui sur $\hat\Sigma\cap\mathbb P(\mathcal M)\simeq X$ induit l'application diagonale. On sait de plus que \deleted{$F_{\mathcal M}(y,x)$} $F_{\mathcal M}$ transforme la symétrie de $X_{\mathcal M}\times X_{\mathcal M}$ en \inserted{l'effet sur $\hat\Sigma_{\mathcal M}$ de} l'anti involution

\[
v\longmapsto v'=\operatorname{tr}(v).1-v
\]

(qui conserve le déterminant \ldots), ce qui pour $v=u+t1$, \inserted{($u\in\Gamma\unsure{\gamma}\mathcal M$, $t\in\Gamma(\mathcal O_S)$)} s'écrit

\[
v'=2t.1-(u+t.1)=-u+t1
\]

soit $-v'=u-t.1$

% archive sheet number: 109
